\documentclass{report}

% Packages for images.
\usepackage{graphicx}
\usepackage{float}

% Packages for mathematical formulae/symbols.
\usepackage{amsmath}
\usepackage{amssymb}
\usepackage{amsfonts}

% Packages for code snippets and algorithms.
\usepackage{listings}
\usepackage{xcolor}
\definecolor{codegreen}{RGB}{0,128,0}
\definecolor{codegray}{RGB}{110,110,110}
\definecolor{codepurple}{RGB}{163,21,21}
\definecolor{codeblue}{RGB}{0,0,180}
\definecolor{backwhite}{RGB}{255,255,255}
\definecolor{backcolour}{RGB}{248,248,248}

\lstdefinestyle{myRstyle}{
    language=R,
    backgroundcolor=\color{backwhite},
    basicstyle=\ttfamily\small,
    keywordstyle=\color{codeblue}\bfseries,
    commentstyle=\color{codegreen}\itshape,
    stringstyle=\color{codepurple},
    numberstyle=\tiny\color{codegray},
    numbers=left,
    numbersep=8pt,
    frame=single,
    rulecolor=\color{black!20},
    breaklines=true,
    showstringspaces=false,
    tabsize=2,
    keepspaces=true,
    otherkeywords={0,1,2,3,4,5,6,7,8,9,<-},
    morekeywords={TRUE,FALSE,NULL,NA,NaN,Inf}
}

\lstset{style=myRstyle}
\usepackage[ruled,vlined]{algorithm2e}

% Creating the Example and Proof boxes.
\definecolor{exampleback}{RGB}{240, 248, 240}
\definecolor{exampleframe}{RGB}{34, 112, 62}
\definecolor{proofback}{RGB}{240, 244, 250}
\definecolor{proofframe}{RGB}{25, 80, 150}
\usepackage[most]{tcolorbox}
\newtcolorbox{examplebox}[1][]{
    colback=exampleback,
    colframe=exampleframe,
    coltitle=white,
    fonttitle=\bfseries,
    title={#1},
    arc=2pt,
    boxrule=1pt
}
\newtcolorbox{proofbox}[1][]{
    colback=proofback,
    colframe=proofframe,
    coltitle=white,
    fonttitle=\bfseries,
    title={Proof: #1},
    arc=2pt,
    boxrule=1pt
}

% Packages for styling and layout.
\usepackage{booktabs}
\usepackage{array}
\usepackage{geometry}
\usepackage{hyperref}

% Packages for plotting graphs.
\usepackage{tikz}
\usepackage{pgfplots}
\pgfplotsset{compat=1.18}

% Packages for compilation.
\usepackage{microtype}

\begin{document}

\begin{titlepage}
    \centering
    \vspace*{2cm}
    
    % Subject Title.
    {\Huge \bfseries Fundamentals of Statistical Learning - Unit 1 \par}
    \vspace{1.5cm}
    
    % Master's Degree and Faculty.
    {\Large Master's Degree in Data Science, Sapienza Università di Roma \par}
    {\large Faculty of Information Engineering, Informatics and Statistics \par}
    \vspace{0.5cm}
    
    % Department.
    {\large Department of Computer Science \par}
    \vspace{3cm}
    
    % Author.
    {\Large \textbf{Author:} Gianmaria Romano \par}
    \vspace{3cm}

    % Reference Bibliography
    {\Large Reference Bibliography \par}
    {\large F. M. Dekking, C. Kraaikamp, H. P. Lopuhaa, L. E. Meester: \href{https://cis.temple.edu/~latecki/Courses/CIS2033-Spring13/Modern_intro_probability_statistics_Dekking05.pdf}{\textit{A Modern Introduction to Probability and Statistics}} \\
    L. Wasserman: \href{https://egrcc.github.io/docs/math/all-of-statistics.pdf}{\textit{All of Statistics: A Concise Course in Statistical Inference}}} \\
    N. Matloff: \href{https://events.psai.ph/2018/ac/downloads/trainings/T3/book_matloff.pdf}{\textit{The Art of R Programming}} \\
    \vfill
    
    % Academic Year.
    {\large Academic Year 2026/2027 \par}
    \vspace*{1cm}
\end{titlepage}

\clearpage

%\maketitle
\begin{tcolorbox}[colback=gray!5!white, colframe=gray!75!black, title=Disclaimer]
    This document contains notes for the \textbf{Fundamentals of Statistical Learning - Unit 1} course, given by \textbf{Professor Brutti} as part of the Master’s Degree in Data Science at \textbf{Sapienza Università di Roma} during the \textbf{first} semester of Academic Year \textbf{2026/2027}.
    
    \smallskip
    \textbf{Note: }These notes may not cover all details of each lecture, so it is highly suggested to consult the Professor's official resources and materials. 
    
    \smallskip
    \textbf{Important: }These notes are free to share and use, but please remember to credit me as the author and do not try to hide/remove this page.
\end{tcolorbox}

\begin{tcolorbox}[colback=gray!5!white, colframe=gray!75!black, title=Contacts]
    For further information or requests, feel free to \href{https://t.me/toritosenso}{contact me privately} or reach out through my repository: \href{https://github.com/GianmariaRomano/Data-Science-Notes}{https://github.com/GianmariaRomano/Data-Science-Notes}.
\end{tcolorbox}

\begin{tcolorbox}[colback=gray!5!white, colframe=gray!75!black, title=License]
    These documents are distributed under the \href{https://www.gnu.org/licenses/gpl-3.0.en.html}{GNU General Public License 3.0}, a form of copyleft intended for use on a manual, textbook or other documents. Material licensed under the current version of the license can be used for any purpose, as long as the use meets certain conditions:
    \begin{itemize}
        \item All previous authors of the work must be attributed.
        \item All changes to the work must be logged.
        \item All derivative works must be licensed under the same license.
        \item The full text of the license, unmodified invariant sections as defined by the author if any, and any other added warranty disclaimers (such as a general disclaimer alerting readers that the document may not be accurate for example) and copyright notices from previous versions must be maintained.
        \item Technical measures such as DRM may not be used to control or obstruct distribution or editing of the document.
    \end{itemize}
\end{tcolorbox}

\tableofcontents

\chapter{Defining probability}
\section{Introducing a qualitative idea of probability}
Probability is the field of mathematics that studies events and how likely they are to occur in order to describe relationships or patterns in existing models. \newline
Starting from a model's data generation process, probability aims to investigate the properties of the observed outcomes, whereas inference makes use of the observed outcomes to study the underlying data generation process.
\begin{examplebox}[\textbf{Example: }Diffusion models]
    Diffusion models are a type of generative model that create new data by gradually reversing a noising process. \newline
    From a mathematical perspective, a diffusion model can be described using a \textit{Markovian Direct Acyclic Graph} whose nodes represent the observable or latent variables involved in the process and whose edges denote the transition probabilities of the model.
\end{examplebox}
\section{Formalizing probability using set theory}
It is possible to formalize probabilistic experiments using set theory. \newline
In fact, starting from a sample space $\displaystyle \Omega$, which represents the set of all possible outcomes, an event $\displaystyle A$ on $\displaystyle \Omega$ can be defined as the set of all favourable outcomes in which the event does occur.
\begin{examplebox}[\textbf{Example: }Tossing a coin]
    Suppose you toss a coin: it is possible to construct a random experiment by defining the sample space $\displaystyle \Omega = \{Heads, Tails\}$ and considering the event $\displaystyle A = \{Heads\}$.
\end{examplebox}
\subsection{Set operations on probability spaces}
Given two events $\displaystyle A$ and $\displaystyle B$ on the same sample space $\displaystyle \Omega$, it is possible to define three main operations:
\begin{itemize}
    \item \textbf{Union: }The set $\displaystyle A \cup B = \{\omega \in \Omega: \omega \in A \text{ or } \omega \in B\}$ corresponds to the event in which at least one of $\displaystyle A$ or $\displaystyle B$, or both, occurs. \newline
    In particular, it is always the case that $\displaystyle A \cup \Omega = \Omega$ and $\displaystyle A \cup \emptyset = A$.
    \item \textbf{Intersection: }The set $\displaystyle A \cap B = \{\omega \in \Omega : \omega \in A \text{ and } \omega \in B\}$ corresponds to the event in which both $\displaystyle A$ and $\displaystyle B$ occur. \newline
    In particular, it is always the case that $\displaystyle A \cap \Omega = A$ and $\displaystyle A \cap \emptyset = \emptyset$.
    \item \textbf{Complement: }The set $\displaystyle A^C = \Omega \setminus A = \{\omega \in \Omega: \omega \notin A\}$ corresponds to the event in which $\displaystyle A$ does not occur. \newline
    In particular, it is always the case that $\displaystyle A \cup A^C = \Omega$ and $\displaystyle A \cap A^C = \emptyset$.
\end{itemize}
Additionally, given a sequence of events $\displaystyle \{A_1, \dots, A_n\}$ on the same sample space $\displaystyle \Omega$, De Morgan's laws allow to switch between the union and intersection operations through the following identities:
\[\Bigg(\bigcup_{i = 1}^n A_i\Bigg)^C = \bigcap_{i = 1}^n A_i^C \text{ and } \Bigg(\bigcap_{i = 1}^n A_i\Bigg)^C = \bigcup_{i = 1}^n A_i^C\]
\subsection{Sample space partitions}
A sequence of events $\displaystyle \{A_1, \dots, A_n\}$ is said to form a partition on a sample space $\displaystyle \Omega$ if the following conditions are satisfied:
\begin{enumerate}
    \item $\displaystyle A_1 \cup \dots \cup A_n = \Omega$, typically requiring $\displaystyle A_i \neq \emptyset \text{ } \forall i =1, \dots, n$ as well.
    \item The sets in the sequence are pairwise disjoint, meaning that $\displaystyle A_i \cap A_j = \emptyset \text{ } \forall i \neq j$.
\end{enumerate}
Generally speaking, partitions come in handy to explain the probability of an event by decomposing it as a union of disjoint sets with respect to a sequence of mutually exclusive hypotheses.
\section{Probability measure functions}
While several definitions of probability exist, Kolmogorov's axiomatic definition is the most commonly used when defining probability measures. \newline
According to this definition, a probability measure function $\displaystyle \mathbb{P}: \mathcal{P}(\Omega) \to [0, 1]$ assigns to each event $\displaystyle A \subseteq \Omega$ a probability $\displaystyle \mathbb{P}(A)$ such that:
\begin{enumerate}
    \item $\displaystyle \mathbb{P}(A) \in [0, 1] \text{ } \forall A \subseteq \Omega$.
    \item $\displaystyle \mathbb{P}(\emptyset) = 0$ and $\displaystyle \mathbb{P}(\Omega) = 1$.
    \item $\displaystyle \mathbb{P}(A \cup B) = \mathbb{P}(A) + \mathbb{P}(B) - \mathbb{P}(A \cap B)$, with $\displaystyle \mathbb{P}(A \cup B) = \mathbb{P}(A) + \mathbb{P}(B)$ if and only if $\displaystyle A \cap B = \emptyset$. \newline
    This last axiom can be generalized for a sequence of $\displaystyle n$ events through the inclusion-exclusion principle, which states that
    \[\mathbb{P}\Bigg(\bigcup_{i = 1}^n A_i\Bigg) = \sum_{k = 1}^n (-1)^{k + 1}\Bigg(\sum_{1 \leq i_1 < \dots < i_k \leq n} \mathbb{P}(A_{i_1} \cap \dots A_{i_k})\Bigg)\]
\end{enumerate}
\begin{examplebox}[\textbf{Example: }Checking polynomial identities using a randomized algorithm]
    Consider two polynomials of degree $\displaystyle d$:
    \[F(x) = \prod_{i = 1}^d (x - a_i) \text{ and } G(x) = \sum_{i = 1}^d c_ix^i\]
    A simple way of checking whether $\displaystyle F(x) \equiv G(x)$ involves expanding $\displaystyle F(x)$ into its canonical form and compare the resulting coefficient with those of $\displaystyle G(x)$, although this algorithm would require performing $\displaystyle \Theta(d^2)$ operations. \newline
    Knowing that, for any randomly chosen point $\displaystyle x_0$, evaluating $\displaystyle F(x_0)$ and $\displaystyle G(x_0)$ requires $\displaystyle O(d)$ time, it is possible to cut down complexity using a randomized algorithm:
    \begin{algorithm}[H]
        \caption{Randomized Algorithm for Polynomial Identities.}
        \KwData{$\displaystyle F(x), G(x)$}
        \KwResult{Whether $\displaystyle F(x) \equiv G(x)$.}

        Generate $\displaystyle r$ uniformly at random from a set $\displaystyle \mathcal{I} = \{1, \dots, 100d\}$\;
        Compute $\displaystyle F(r)$ and $\displaystyle G(r)$\;
        \eIf{$\displaystyle F(r) = G(r)$}{
            $\displaystyle F(x) \equiv G(x)$\;
        }{
            $\displaystyle F(x) \not\equiv G(x)$\;
        }
    \end{algorithm}
    \noindent However, the randomized nature of this algorithm can result in a wrong evaluation. \newline
    In fact, while $\displaystyle F(x) \equiv G(x)$ results in $\displaystyle F(r) = G(r) \text{ } \forall r$, it might be the case that, for some $\displaystyle r$, $\displaystyle F(r) = G(r)$ even if $\displaystyle F(x) \not\equiv G(x)$: this implies that the algorithm has one-sided error as it can make mistakes only when deciding that $\displaystyle F(x) \equiv G(x)$. \newline
    Therefore, define the event
    \[E = \{F(r) = G(r) \text{ but } F(x) \not\equiv G(x)\}\]
    This event occurs if and only if $\displaystyle r$ is a root of the polynomial $\displaystyle H(x) = F(x) - G(x)$: since $\displaystyle H(x)$ will always have degree at most $\displaystyle d$, it will have at most $\displaystyle d$ roots, allowing to bound the error probability as follows:
    \[\mathbb{P}(E) = \frac{|\{r \in \mathcal{I} : r \text{ is a root of } H(x)\}|}{|\mathcal{I}|} \leq \frac{d}{100d} = \frac{1}{100}\]
\end{examplebox}
\noindent As a consequence to Kolmogorov's axiomatic definition of probability, every probability space $\displaystyle (\Omega, \mathbb{P})$ satisfies the property that $\displaystyle \mathbb{P}(A^C) = 1 - \mathbb{P}(A) \text{ } \forall A \subseteq \Omega$.
\begin{proofbox}
    Consider the partition $\displaystyle \Omega = A \cup A^C$ and apply Kolmogorov's axioms.
    \[1 = \mathbb{P}(\Omega) = \mathbb{P}(A \cup A^C) = \mathbb{P}(A) + \mathbb{P}(A^C) \Rightarrow \mathbb{P}(A^C) = 1 - \mathbb{P}(A)\]
\end{proofbox}
\subsection{Probability bounds}
Boole's inequality states that, for every sequence of events $\displaystyle \{A_1, \dots, A_n\}$, it holds that
\[\mathbb{P}\Bigg(\bigcup_{i = 1}^n A_i\Bigg) \leq \sum_{i = 1}^n \mathbb{P}(A_i)\]
\begin{proofbox}
    This inequality can be proved by induction on the length of the sequence.
    \begin{itemize}
        \item For the base case, give by $\displaystyle n = 1$, the inequality is trivially true as $\displaystyle \mathbb{P}(A_1) \leq \mathbb{P}(A_1)$.
        \item Assuming the inequality is true for some $\displaystyle n$, the goal is to show that it is true for $\displaystyle (n + 1)$ as well. \newline
        Reformulate $\displaystyle A_1 \cup \dots \cup A_{n + 1} = (A_1 \cup \dots \cup A_n) \cup A_{n + 1}$ and let
        \begin{align*}
            \mathbb{P}\Bigg(\bigcup_{i = 1}^{n + 1} A_i\Bigg) 
              &= \mathbb{P}\Bigg(\Bigg(\bigcup_{i = 1}^n A_i\Bigg) \cup A_{n + 1}\Bigg) = \mathbb{P}\Bigg(\bigcup_{i = 1}^n A_i\Bigg) + \mathbb{P}(A_{n + 1}) - \mathbb{P}\Bigg(\Bigg(\bigcup_{i = 1}^n A_i\Bigg) \cap A_{n + 1}\Bigg) \\
              &\stackrel{\mathbb{P}(.) \ge 0}{\le} \mathbb{P}\Bigg(\bigcup_{i = 1}^n A_i\Bigg) + \mathbb{P}(A_{n + 1}) \stackrel{\text{hp. ind.}}{\le} \sum_{i=1}^{n} \mathbb{P}(A_i) + \mathbb{P}(A_{n+1}) = \sum_{i=1}^{n+1} \mathbb{P}(A_i)
        \end{align*}
    \end{itemize}
    Having proved the base case and the inductive case, this inequality must be true $\displaystyle \forall n \in \mathbb{N}$, which implies that probability measure functions are always sub-additive.
\end{proofbox}
\noindent As a consequence to Boole's inequality, Bonferroni's inequality states that, for every sequence of events $\displaystyle \{A_1, \dots, A_n\}$, it must be the case that
\[\mathbb{P}\Bigg(\bigcap_{i = 1}^n A_i\Bigg) \ge \sum_{i = 1}^n \mathbb{P}(A_i) - (n - 1)\]
\begin{proofbox}
    Reformulate the joint probability of the events using De Mrogan's laws.
    \begin{align*}
        \mathbb{P}\Bigg(\bigcap_{i = 1}^n A_i\Bigg)
            &= 1 - \mathbb{P}\Bigg(\Bigg(\bigcap_{i = 1}^n A_i\Bigg)^C\Bigg) = \stackrel{\text{De Morgan}}{=} 1 - \mathbb{P}\Bigg(\bigcup_{i = 1}^n A_i^C\Bigg) \\
            &\stackrel{\text{Boole}}{\ge} 1 - \sum_{i = 1}^n \mathbb{P}(A_i^C) = 1 - \sum_{i = 1}^n (1 - \mathbb{P}(A_i)) \\
            &= 1 - \Bigg(\sum_{i = 1}^n 1 - \sum_{i = 1}^n \mathbb{P}(A_i)\Bigg) = \sum_{i = 1}^n \mathbb{P}(A_i) - (n - 1)
    \end{align*}
\end{proofbox}

\chapter{Conditional probability}
\section{Understanding conditional probability}
Given a sample space $\displaystyle \Omega$, the conditional probability that event $\displaystyle A$ occurs knowing that event $\displaystyle B$ was observed can be defined as follows:
\[\mathbb{P}(A|B) = \frac{\mathbb{P}(A \cap B)}{\mathbb{P}(B)} \text{, with $\mathbb{P}(B) > 0$.}\]
\begin{examplebox}[\textbf{Example: }Measles and red spots]
    Consider the two events
    \[A = \{\text{I have measles}\} \text{ and } B = \{\text{I have red spots on my face}\}\]
    Knowing that red spots are a common symptom of measles, it will be the case that $\displaystyle \mathbb{P}(B|A) = 1$. \newline
    However, the same thing cannot be said about $\displaystyle \mathbb{P}(A|B)$ as red spots might be caused by other external factors, such as eating hot food or getting a scratch.
\end{examplebox}
\noindent From a broader perspective, a conditional probability distribution can be described using a probability measure function that satisfies the following properties:
\begin{enumerate}
    \item $\displaystyle \mathbb{P}(A|B) \in [0, 1] \text{ } \forall A \subseteq \Omega$.
    \item $\displaystyle \mathbb{P}(\Omega|B) = 1$.
    \item For any sequence of disjoint events $\displaystyle \{A_1, \dots A_n\}$, it is always the case that
    \[\mathbb{P}\Bigg(\bigcup_{i = 1}^n A_i|B\Bigg) = \sum_{i = 1}^n \mathbb{P}(A_i|B)\]
    As a consequence, it is always the case that $\displaystyle \mathbb{P}(A^C|B) = 1 - \mathbb{P}(A|B) \text{ } \forall A \subseteq \Omega$.
    \begin{proofbox}
        Let $\displaystyle B = B \cap \Omega$ on the partition $\displaystyle \Omega = A \cup A^C$.
        \begin{align*}
            \mathbb{P}(B)
                &= \mathbb{P}(B \cap \Omega) = \mathbb{P}(B \cap (A \cup A^C)) = \mathbb{P}((B \cap A) \cup (B \cap A^C)) \\
                &= \mathbb{P}(A \cap B) + \mathbb{P}(A^C \cap B) \Rightarrow \mathbb{P}(A^C \cap B) = \mathbb{P}(B) - \mathbb{P}(A \cap B)
        \end{align*}
        Therefore, applying the definition of conditional probability will result in
        \[\mathbb{P}(A^C|B) = \frac{\mathbb{P}(A^C \cap B)}{\mathbb{P}(B)} = \frac{\mathbb{P}(B) - \mathbb{P}(A \cap B)}{\mathbb{P}(B)} = \frac{\mathbb{P}(B)}{\mathbb{P}(B)} - \frac{\mathbb{P}(A \cap B)}{\mathbb{P}(B)} = 1 - \mathbb{P}(A|B)\]
    \end{proofbox}
\end{enumerate}
\subsection{The (conditional) odds of an event}
Generally speaking, the odds of an event can be seen as the likelihood that the event will occur relative to it not occurring. \newline
For this reason, the odds of an event $\displaystyle A \subseteq \Omega$ with probability $\displaystyle \mathbb{P}(A) \in [0, 1)$ can be computed as
\[\mathcal{O}(A) = \frac{\mathbb{P}(A)}{\mathbb{P}(A^C)} = \frac{\mathbb{P}(A)}{1 - \mathbb{P}(A)}\]
More specifically, the conditional odds of event $\displaystyle A$ given that a second event $\displaystyle B$ occurred can be computed as
\[\mathcal{O}(A|B) = \frac{\mathbb{P}(A|B)}{\mathbb{P}(A^C|B)} = \frac{\mathbb{P}(A \cap B)}{\mathbb{P}(A^C \cap B)}\]
\subsection{The multiplication rule}
The multiplication rule allows to reformulate the joint probability of two events $\displaystyle A$ and $\displaystyle B$ according to the following equation:
\[\mathbb{P}(A \cap B) = \mathbb{P}(A|B)\mathbb{P}(B) = \mathbb{P}(B|A)\mathbb{P}(A)\]
In particular, repeatedly applying the multiplication rule to a sequence of events $\displaystyle \{A_1, \dots, A_n\}$ results in the chain rule, which can be formulated as
\[\mathbb{P}\Bigg(\bigcap_{i = 1}^n A_i\Bigg) = \mathbb{P}(A_1)\cdot\mathbb{P}(A_2|A_1)\cdot\dots\cdot\mathbb{P}(A_n|A_1 \cap \dots A_{n + 1})\]
\section{Defining event independence}
Two events $\displaystyle A$ and $\displaystyle B$ on a sample space $\displaystyle \Omega$ are said to be independent if knowing that $\displaystyle B$ occurred does not add information about the probability of $\displaystyle A$ occurring, resulting in
\[A \bot B \Leftrightarrow \mathbb{P}(A|B) = \mathbb{P}(A) \text{ and } \mathbb{P}(B|A) = \mathbb{P}(B) \text{, resulting in $\mathbb{P}(A \cap B) = \mathbb{P}(A)\mathbb{P}(B)$.}\]
Depending on the context, event independence can be assumed or verified, although it is always possible to show that, if $\displaystyle A$ and $\displaystyle B$ are events such that $\displaystyle \mathbb{P}(A), \mathbb{P}(B) > 0$ and $\displaystyle A \cap B = \emptyset$, then they events cannot be independent.
\begin{proofbox}
    It is possible to notice that having $\displaystyle A \cap B = \emptyset$ and knowing that one of the two events occurred would violate the definition of event independence because
    \[\mathbb{P}(A)\mathbb{P}(B) > 0 \text{, but $\mathbb{P}(A \cap B) = \mathbb{P}(\emptyset) = 0$.}\]
\end{proofbox}
\section{Generalizing event independence}
For $\displaystyle n>2$ events, independence can be defined at several distinct levels.
\subsection{Pairwise and mutual independence}
A (countable) sequence of events $\displaystyle \{A_1, \dots, A_n\}$ is said to be pairwise independent if
\[A_i \bot A_j \text{ } \forall i \neq j \Leftrightarrow \mathbb{P}(A_i \cap A_j) = \mathbb{P}(A_i)\mathbb{P}(A_j) \text{ } \forall i \neq j\]
Additionally, a (countable) sequence of events $\displaystyle \{A_1, \dots, A_n\}$ is said to be mutually independent if, for any group of $\displaystyle k$ events in the sequence, it is the case that
\[\mathbb{P}\Bigg(\bigcap_{i = 1}^k A_{n_i}\Bigg) = \prod_{i = 1}^k \mathbb{P}(A_{n_i}) \text{, for $1 \leq k \leq n$ and $k \in \mathbb{N}$.}\]
It should be noticed that, while mutual independence always implies pairwise independence, the converse statement is not true because a sequence of events can be pairwise independent but not mutually independent.
\begin{examplebox}[\textbf{Example: }A special die]
    Given a balanced four-faced die, which has a red face, a green face, a blue face and a striped face with all three colours, define the events
    \[R = \{\text{I see red}\} \text{, } G = \{\text{I see green}\} \text{ and } B = \{\text{I see blue}\}\]
    It is possible to show that these events are pairwise independent as
    \begin{align*}
        \mathbb{P}(R \cap G) &= 0.25 = 0.5 \cdot 0.5 = \mathbb{P}(R)\mathbb{P}(G) \\
        \mathbb{P}(R \cap B) &= 0.25 = 0.5 \cdot 0.5 = \mathbb{P}(R)\mathbb{P}(B) \\
        \mathbb{P}(G \cap B) &= 0.25 = 0.5 \cdot 0.5 = \mathbb{P}(G)\mathbb{P}(B)
    \end{align*}
    However, the events are not mutually independent because
    \[\mathbb{P}(R \cap G \cap B) = 0.25 \neq 0.125 = 0.5 \cdot 0.5 \cdot 0.5 = \mathbb{P}(R)\mathbb{P}(G)\mathbb{P}(B)\]
\end{examplebox}
\subsection{Conditional independence}
Two events $\displaystyle A$ and $\displaystyle B$ are said to be conditionally independent given an event $\displaystyle C$ if, knowing that $\displaystyle C$ occurred, observing event $\displaystyle B$ does not add information about the probability of $\displaystyle A$ occurring, resulting in
\[\mathbb{P}(A \cap B|C) = \mathbb{P}(A|C)\mathbb{P}(B|C) \text{, with $\mathbb{P}(C) > 0$.}\]
\begin{examplebox}[\textbf{Example: }University admissions]
    For a student applying to a Master's Degree course, define the events
    \[DS = \{\text{Accepted to Data Science}\} \text{ and } ST = \{\text{Accepted to Statistics}\}\]
    Knowing that event $\displaystyle DS$ occurred can provide additional information about the likelihood of event $\displaystyle ST$ occurring, meaning that the two events are not independent. \newline
    However, if the admission procedure is solely based on the student's GPA $\displaystyle C$, knowing that one event occurred does not provide additional information about the other event occurring as well, making $\displaystyle DS$ and $\displaystyle ST$ conditionally independent given event $\displaystyle C$.
\end{examplebox}
\section{The law of total probability}
Given a sample space $\displaystyle \Omega$, the law of total probability allows to explain the probability of an event $\displaystyle A \subseteq \Omega$ by decomposing it over the partition $\displaystyle \Omega = B \cup B^C$, resulting in
\[\mathbb{P}(A) = \mathbb{P}(A \cap B) + \mathbb{P}(A \cap B^C) = \mathbb{P}(A|B)\mathbb{P}(B) + \mathbb{P}(A|B^C)\mathbb{P}(B^C)\]
\begin{proofbox}
    Let $\displaystyle A = A \cap \Omega$ on the partition $\displaystyle \Omega = B \cup B^C$.
    \begin{align*}
        \mathbb{P}(A)
            &= \mathbb{P}(A \cap \Omega) = \mathbb{P}(A \cap (B \cup B^C)) = \mathbb{P}((A \cap B) \cup (A \cap B^C)) \\
            &= \mathbb{P}(A \cap B) + \mathbb{P}(A \cap B^C) \stackrel{\text{Multiplication rule}}{=} \mathbb{P}(A|B)\mathbb{P}(B) + \mathbb{P}(A|B^C)\mathbb{P}(B)
    \end{align*}
\end{proofbox}
\noindent Additionally, it is possible to generalize the law of total probability to explain the probability of an event $\displaystyle E$ with respect to a sequence of mutually exclusive hypotheses $\displaystyle \{H_1, \dots, H_n\}$ as
\[\mathbb{P}(E) = \sum_{i = 1}^n \mathbb{P}(E|H_i) = \sum_{i = 1}^n \mathbb{P}(E|H_i)\mathbb{P}(H_i)\]
\section{Bayes' theorem}
Bayes' theorem allows to switch between conditional probabilities by setting
\[\mathbb{P}(A|B) = \frac{\mathbb{P}(A \cap B)}{\mathbb{P}(B)} = \frac{\mathbb{P}(B|A)\mathbb{P}(A)}{\mathbb{P}(B)} \text{ with $\mathbb{P}(B) > 0$.}\]
In particular, when considering a sequence of mutually exclusive hypotheses $\displaystyle \{H_1, \dots, H_n\}$ with known prior probabilities, Bayes' theorem allows to compute the posterior probability of hypothesis $\displaystyle H_i$ given the observed evidence $\displaystyle E$ as
\[\mathbb{P}(H_i|E) = \frac{\mathbb{P}(E|H_i)\mathbb{P}(H_i)}{\mathbb{P}(E)} = \frac{\mathbb{P}(E|H_i)\mathbb{P}(H_i)}{\sum_{j = 1}^n \mathbb{P}(E|H_j)\mathbb{P}(H_j)}\]
In fact, the posterior probability of each hypothesis given the observed evidence will be computed in terms of the likelihood $\displaystyle \mathbb{P}(E|H_i)$, which describes the probability of seeing the evidence given that the hypothesis $\displaystyle H_i$ is true, the prior probability $\displaystyle \mathbb{P}(H_i)$, which describes the probability of $\displaystyle H_i$ being true, and the marginal probability $\displaystyle \mathbb{P}(E)$, which denotes the total probability of seeing the evidence under all possible hypotheses.
\begin{examplebox}[\textbf{Example: }Iterated Bayes for medical diagnosis]
    Consider the events
    \[T = \{\text{I test positive to a disease}\} \text{ and } D = \{\text{I actually have the disease}\}\]
    In particular, the disease affects $\displaystyle 0.1\%$ of the population and it is known that a dedicated test correctly detects $\displaystyle 99\%$ of disease case but produces a false positive result on $\displaystyle 1\%$ of occasions, which means that
    \[\mathbb{P}(D|P) = \frac{\mathbb{P}(P|D)\mathbb{P}(D)}{\mathbb{P}(P|D)\mathbb{P}(D) + \mathbb{P}(P|D^C)\mathbb{P}(D^C)} = \frac{0.99 \cdot 0.001}{(0.99 \cdot 0.001) + (0.01 \cdot 0.999)} = 9\%\]
    Since the posterior probability appears to be quite low, Bayes' rule is applied once again, using the posterior probability as the new prior probability, resulting in
    \[\mathbb{P}(D|P) = \frac{\mathbb{P}(P|D)\mathbb{P}(D)}{\mathbb{P}(P|D)\mathbb{P}(D) + \mathbb{P}(P|D^C)\mathbb{P}(D^C)} = \frac{0.99 \cdot 0.09}{(0.99 \cdot 0.09) + (0.01 \cdot 0.91)} = 91\%\]
\end{examplebox}
\subsection{The posterior odds of a hypothesis}
When considering an event $\displaystyle E$ on a sequence of mutually exclusive hypotheses $\displaystyle \{H_1, \dots, H_n\}$, it is possible to compute the posterior odds of hypothesis $\displaystyle H_i$ given evidence $\displaystyle E$ as
\[\mathcal{O}(H_i|E) = \frac{\mathbb{P}(H_i|E)}{\mathbb{P}(H_i^C|E)} = \frac{\frac{\mathbb{P}(E|H_i)\mathbb{P}(H_i)}{\mathbb{P}(E)}}{\frac{\mathbb{P}(E|H_i^C)\mathbb{P}(H_i^C)}{\mathbb{P}(E)}} = \frac{\mathbb{P}(E|H_i)\mathcal{O}(H_i)}{\mathbb{P}(E|H_i)}\]

\end{document}